\sectionheading{5.19 কম্পিউটেশনাল থিংকিং (Computational Thinking)}

কোনো জটিল Problem দেখেই সরাসরি Code লেখা শুরু করা সবসময়
কার্যকর নয়।

প্রথমে Problem-টিকে বুঝতে হয়, ছোট অংশে ভাগ করতে হয়,
প্রয়োজনীয় তথ্য আলাদা করতে হয় এবং ধাপে ধাপে Solution
তৈরি করতে হয়।

সমস্যা সমাধানের এই যৌক্তিক ও সুশৃঙ্খল চিন্তার পদ্ধতিই
\textbf{Computational Thinking}।

%================================================
% Definition
%================================================
\booksubheading{Computational Thinking কাকে বলে?}

\begin{definitionbox}

\textbf{\color{darkTeal}সংজ্ঞা}

কোনো Problem-কে এমনভাবে
\textbf{বিশ্লেষণ, বিভাজন ও ধাপে ধাপে সমাধানের পরিকল্পনা}
করার চিন্তাপদ্ধতি, যাতে Solution মানুষ বা Computer
কার্যকরভাবে বাস্তবায়ন করতে পারে, তাকে
\textbf{Computational Thinking} বলে।

\end{definitionbox}

\textbf{সহজভাবে মনে রাখবে—}

\begin{center}

\textbf{
Problem বুঝি
\(\longrightarrow\)
ছোট অংশে ভাগ করি
\(\longrightarrow\)
প্রয়োজনীয় বিষয় খুঁজি
\(\longrightarrow\)
সমাধানের ধাপ তৈরি করি
}

\end{center}

%================================================
% Main Elements
%================================================
\booksubheading{Computational Thinking-এর মূল উপাদান}

Computational Thinking বোঝার জন্য চারটি গুরুত্বপূর্ণ ধারণা হলো—

\begin{center}

\textbf{
Decomposition
\(\longrightarrow\)
Pattern Recognition
\(\longrightarrow\)
Abstraction
\(\longrightarrow\)
Algorithmic Thinking
}

\end{center}

%================================================
% Decomposition
%================================================
\booksubheading{১. Decomposition — সমস্যা বিভাজন}

\begin{definitionbox}

একটি বড় ও জটিল Problem-কে কয়েকটি
\textbf{ছোট ও সহজ Problem-এ ভাগ করার প্রক্রিয়াকে}
\textbf{Decomposition} বলে।

\end{definitionbox}

\begin{examplebox}{সহজ উদাহরণ}

ধরা যাক একটি Student Result System তৈরি করতে হবে।

পুরো Problem-টিকে এভাবে ভাগ করা যায়—

\begin{itemize}

\item শিক্ষার্থীর তথ্য গ্রহণ
\item নম্বর গ্রহণ
\item মোট নম্বর নির্ণয়
\item Average নির্ণয়
\item Grade নির্ণয়
\item Result প্রদর্শন

\end{itemize}

অর্থাৎ একটি বড় Problem কয়েকটি ছোট কাজে বিভক্ত হলো।

\end{examplebox}

\begin{center}

\textbf{
Decomposition
=
বড় Problem
\(\longrightarrow\)
ছোট Problem
}

\end{center}

%================================================
% Pattern Recognition
%================================================
\booksubheading{২. Pattern Recognition — সাদৃশ্য বা Pattern শনাক্তকরণ}

\begin{definitionbox}

Problem-এর বিভিন্ন অংশ বা পূর্বে সমাধান করা Problem-এর মধ্যে
\textbf{মিল, পুনরাবৃত্তি বা একই ধরনের বৈশিষ্ট্য খুঁজে বের করাকে}
\textbf{Pattern Recognition} বলে।

\end{definitionbox}

\begin{examplebox}{উদাহরণ}

একাধিক শিক্ষার্থীর Result তৈরির ক্ষেত্রে প্রত্যেকের জন্যই—

\begin{center}

\textbf{
Marks
\(\longrightarrow\)
Total
\(\longrightarrow\)
Average
\(\longrightarrow\)
Grade
}

\end{center}

একই ধরনের কাজ পুনরাবৃত্তি হচ্ছে।

এই মিল বা পুনরাবৃত্তি শনাক্ত করাই
\textbf{Pattern Recognition}।

\end{examplebox}

Pattern শনাক্ত করতে পারলে আগে জানা Solution নতুন বা একই ধরনের
Problem সমাধানে কাজে লাগানো যায়।

%================================================
% Abstraction
%================================================
\booksubheading{৩. Abstraction — প্রয়োজনীয় তথ্য নির্বাচন}

\begin{definitionbox}

Problem সমাধানের জন্য প্রয়োজনীয় তথ্য রেখে
\textbf{অপ্রয়োজনীয় বা কম গুরুত্বপূর্ণ বিস্তারিত বিষয় উপেক্ষা করার
প্রক্রিয়াকে} \textbf{Abstraction} বলে।

\end{definitionbox}

\begin{examplebox}{সহজ উদাহরণ}

একজন শিক্ষার্থীর Average নির্ণয়ের জন্য প্রয়োজন—

\begin{center}

\textbf{শিক্ষার্থীর প্রাপ্ত নম্বর}
\end{center}

কিন্তু তার প্রিয় রং বা বাসার রং Average নির্ণয়ের জন্য
প্রয়োজনীয় নয়।

অতএব Problem সমাধানের জন্য প্রয়োজনীয় তথ্য নেওয়া হবে এবং
অপ্রাসঙ্গিক তথ্য উপেক্ষা করা হবে।

\end{examplebox}

\begin{center}

\textbf{
প্রয়োজনীয় তথ্য রাখি
+
অপ্রাসঙ্গিক তথ্য বাদ দিই
=
Abstraction
}

\end{center}

%================================================
% Algorithmic Thinking
%================================================
\booksubheading{৪. Algorithmic Thinking — ধাপভিত্তিক সমাধান}

\begin{definitionbox}

Problem সমাধানের জন্য
\textbf{সুস্পষ্ট, যৌক্তিক ও ধারাবাহিক ধাপ নির্ধারণ করার
চিন্তাপদ্ধতিকে}
\textbf{Algorithmic Thinking} বলে।

\end{definitionbox}

\begin{examplebox}{উদাহরণ}

দুটি সংখ্যার বড়টি নির্ণয়ের জন্য—

\begin{enumerate}

\item দুটি সংখ্যা গ্রহণ করি।
\item সংখ্যা দুটি তুলনা করি।
\item বড় সংখ্যাটি নির্ধারণ করি।
\item Result প্রদর্শন করি।

\end{enumerate}

অর্থাৎ Problem-এর Solution-কে ধাপে ধাপে সাজানো হয়েছে।

\end{examplebox}

Algorithm পরবর্তী Topic-এ বিস্তারিত আলোচনা করা হবে।

%================================================
% Combined Example
%================================================
\booksubheading{চারটি ধারণা একসঙ্গে}

ধরা যাক একটি School Result Processing System তৈরি করতে হবে।

\begin{tabularx}{\linewidth}{|>{\bfseries}p{0.30\linewidth}|Y|}
\hline

\rowcolor{softBlue}
\multicolumn{1}{|c|}{\textbf{Computational Thinking-এর অংশ}}
&
\multicolumn{1}{c|}{\textbf{কী করা হবে}}
\\
\hline

Decomposition
&
System-কে Input, Calculation, Grade, Output ইত্যাদি
ছোট অংশে ভাগ করা।
\\
\hline

Pattern Recognition
&
সব শিক্ষার্থীর Result তৈরিতে একই ধরনের কাজ শনাক্ত করা।
\\
\hline

Abstraction
&
Result তৈরির জন্য শুধু প্রয়োজনীয় Data নির্বাচন করা।
\\
\hline

Algorithmic Thinking
&
কাজগুলো কোন ক্রমে হবে তা ধাপে ধাপে নির্ধারণ করা।
\\
\hline

\end{tabularx}

\begin{rememberbox}

\textbf{\color{deepOrange}গুরুত্বপূর্ণ}

Computational Thinking কোনো নির্দিষ্ট Programming Language নয়;
এটি \textbf{Problem Solving-এর একটি চিন্তাপদ্ধতি}।

\end{rememberbox}

%================================================
% Importance
%================================================
\booksubheading{Computational Thinking কেন গুরুত্বপূর্ণ?}

\begin{compacttwocolpoints}

\item জটিল Problem সহজে বোঝা যায়।

\item বড় Problem ছোট অংশে ভাগ করা যায়।

\item একই ধরনের Problem-এর Pattern শনাক্ত করা যায়।

\item অপ্রয়োজনীয় তথ্য বাদ দিয়ে মূল Problem-এ মনোযোগ দেওয়া যায়।

\item Algorithm ও Program তৈরি সহজ হয়।

\item Error কমানো এবং Solution যাচাই করা সহজ হয়।

\item সুশৃঙ্খল ও কার্যকর Solution তৈরি করা যায়।

\end{compacttwocolpoints}

%================================================
% Computational Thinking vs Programming
%================================================
\booksubheading{Computational Thinking ও Programming-এর সম্পর্ক}

Computational Thinking এবং Programming পরস্পরের সঙ্গে সম্পর্কিত,
তবে দুটি একই বিষয় নয়।

\textbf{Computational Thinking} হলো Problem কীভাবে সমাধান করা হবে
তা নিয়ে চিন্তা ও পরিকল্পনা।

\textbf{Programming} হলো সেই Solution-কে Programming Language-এর
মাধ্যমে Computer-এ বাস্তবায়ন করা।

\begin{tabularx}{\linewidth}{|>{\bfseries}p{0.22\linewidth}|Y|Y|}
\hline

\rowcolor{softBlue}
\multicolumn{1}{|c|}{\textbf{বিষয়}}
&
\multicolumn{1}{c|}{\textbf{Computational Thinking}}
&
\multicolumn{1}{c|}{\textbf{Programming}}
\\
\hline

মূল কাজ
&
Problem বিশ্লেষণ ও Solution পরিকল্পনা
&
Solution-কে Code-এ রূপ দেওয়া
\\
\hline

পর্যায়
&
Code লেখার আগেই শুরু হতে পারে
&
সাধারণত Solution নির্ধারণের পর
\\
\hline

Language নির্ভরতা
&
নির্দিষ্ট Programming Language-এর ওপর নির্ভরশীল নয়
&
Programming Language প্রয়োজন
\\
\hline

উদাহরণ
&
Problem ভাগ করা, Pattern খোঁজা
&
Python-এ Program লেখা
\\
\hline

\end{tabularx}

\textbf{সহজভাবে—}

\begin{center}

\textbf{
Computational Thinking
=
কীভাবে সমাধান করব?
}

\vspace{0.5mm}

\textbf{
Programming
=
সেই সমাধান Computer-কে কীভাবে করাব?
}

\end{center}

%================================================
% Relation to Modular Programming
%================================================
\booksubheading{Modular Programming-এর সঙ্গে সম্পর্ক}

আগের Topic-এ Modular Programming-এ বড় Program-কে ছোট
Module-এ ভাগ করার কথা বলা হয়েছে।

Computational Thinking-এর \textbf{Decomposition}-এও
বড় Problem-কে ছোট অংশে ভাগ করা হয়।

তবে দুটি একই বিষয় নয়।

\begin{tabularx}{\linewidth}{|>{\bfseries}p{0.28\linewidth}|Y|}
\hline

\rowcolor{softBlue}
\multicolumn{1}{|c|}{\textbf{বিষয়}}
&
\multicolumn{1}{c|}{\textbf{মূল ধারণা}}
\\
\hline

Decomposition
&
Problem-কে ছোট অংশে ভাঙার চিন্তাপদ্ধতি।
\\
\hline

Modular Programming
&
Program-কে ছোট Module-এ সংগঠিত করার Programming পদ্ধতি।
\\
\hline

\end{tabularx}

%================================================
% Quick Comparison
%================================================
\booksubheading{চারটি উপাদান দ্রুত মনে রাখো}

\begin{tabularx}{\linewidth}{|>{\bfseries}p{0.30\linewidth}|Y|}
\hline

\rowcolor{softBlue}
\multicolumn{1}{|c|}{\textbf{ধারণা}}
&
\multicolumn{1}{c|}{\textbf{Shortcut}}
\\
\hline

Decomposition
&
Divide — বড় Problem-কে ছোট অংশে ভাগ করা
\\
\hline

Pattern Recognition
&
Find Pattern — মিল বা পুনরাবৃত্তি খোঁজা
\\
\hline

Abstraction
&
Focus — প্রয়োজনীয় বিষয়ের ওপর গুরুত্ব দেওয়া
\\
\hline

Algorithmic Thinking
&
Make Steps — সমাধানের ধাপ তৈরি করা
\\
\hline

\end{tabularx}

%================================================
% Final Reminder
%================================================
\begin{rememberbox}

\textbf{\color{deepOrange}মনে রাখবে}

\begin{itemize}

\item \textbf{Computational Thinking} একটি Problem Solving
চিন্তাপদ্ধতি; এটি কোনো Programming Language নয়।

\item এর গুরুত্বপূর্ণ চারটি ধারণা—
\textbf{Decomposition, Pattern Recognition, Abstraction এবং
Algorithmic Thinking}।

\item \textbf{Decomposition} =
বড় Problem-কে ছোট অংশে ভাগ করা।

\item \textbf{Pattern Recognition} =
Problem-এর মধ্যে মিল বা পুনরাবৃত্তি খুঁজে বের করা।

\item \textbf{Abstraction} =
প্রয়োজনীয় তথ্য রাখা এবং অপ্রয়োজনীয় তথ্য বাদ দেওয়া।

\item \textbf{Algorithmic Thinking} =
Solution-এর ধাপগুলো যৌক্তিকভাবে সাজানো।

\item Computational Thinking Program লেখার আগেই
Problem বুঝতে ও Solution পরিকল্পনা করতে সাহায্য করে।

\item দ্রুত মনে রাখার ক্রম—

\begin{center}

\textbf{
Divide
\(\longrightarrow\)
Find Pattern
\(\longrightarrow\)
Focus
\(\longrightarrow\)
Make Steps
}

\end{center}

\end{itemize}

\end{rememberbox}

\small

%================================================
% 5.20 Problem Analysis and IPO Model
%================================================
